\section{Preparación del entorno de trabajo}

En esta sección se describen los requisitos hardware y software necesarios para el desarrollo del proyecto, así como la configuración del entorno de trabajo adoptado.

La aplicación se ha desarrollado siguiendo una arquitectura full-stack monolítica basada en Next.js, donde el frontend y la lógica de servidor (API Routes) conviven dentro del mismo proyecto. Para la persistencia de datos se emplea PostgreSQL, ejecutado como servicio independiente mediante contenedores Docker.

Asimismo, se documenta la configuración del sistema de control de versiones y el entorno de documentación utilizado para la redacción de esta memoria.

\subsection{Requisitos hardware}

El desarrollo del proyecto se realizó en dos equipos con las siguientes características:

\begin{itemize}
    \item Sistema operativo: Windows 11 64 bits
    \item Procesador: Intel Core i5 8600K (equipo de sobremesa) / Intel Core i5 11400H (portátil)
    \item Memoria RAM: 16 GB
    \item Espacio en disco disponible: superior a 10 GB
\end{itemize}

No se requieren características hardware especiales más allá de las necesarias para ejecutar herramientas modernas de desarrollo web y contenedores Docker para la virtualización de servicios.

\subsection{Requisitos software}

Las herramientas utilizadas durante el desarrollo del proyecto han sido las siguientes:

\subsubsection*{Entorno de desarrollo}

\begin{itemize}
    \item Node.js v24.13.1
    \item npm 11.10.1 (incluido con Node.js)
    \item Next.js 16.1.6
    \item React 19.x
    \item PostgreSQL 16 (ejecutado en contenedor Docker)
    \item Docker Desktop 4.x
    \item Paquete \texttt{pg} (cliente oficial de PostgreSQL para Node.js)
\end{itemize}

Next.js permite integrar en un único proyecto tanto la interfaz de usuario como la lógica de servidor mediante su sistema de rutas y API interna. Esto elimina la necesidad de un backend independiente basado en frameworks como Express, simplificando la arquitectura general del sistema.

\subsubsection*{Herramientas adicionales}

\begin{itemize}
    \item Git 2.34.0.windows.1
    \item GitHub Desktop 3.5.4
    \item Visual Studio Code 1.109.5
    \item MiKTeX 26.1 (distribución \LaTeX)
    \item Strawberry Perl v5.42.0 (necesario para la ejecución de \texttt{latexmk})
\end{itemize}

\subsection{Instalación del entorno local}

\subsubsection*{Node.js y npm}

Se instaló la versión LTS de Node.js desde su sitio oficial.  
Tras la instalación, se verificó su correcta integración en el sistema mediante:

\begin{lstlisting}[language=bash]
node -v
npm -v
\end{lstlisting}

\subsubsection*{Creación del proyecto Next.js}

El proyecto se inicializó mediante el comando:

\begin{lstlisting}[language=bash]
npx create-next-app@latest nombre-proyecto
\end{lstlisting}

Este comando genera automáticamente la estructura base del proyecto e instala las dependencias principales:

\begin{itemize}
    \item next
    \item react
    \item react-dom
\end{itemize}

Asimismo, configura los scripts de desarrollo y compilación en el archivo \texttt{package.json}.

Para iniciar el servidor de desarrollo se ejecuta:

\begin{lstlisting}[language=bash]
npm run dev
\end{lstlisting}

El servidor queda accesible en:

\begin{center}
\texttt{http://localhost:3000}
\end{center}

\subsubsection*{Configuración de la base de datos PostgreSQL}

La base de datos PostgreSQL se ejecuta mediante un contenedor Docker definido en el archivo \texttt{docker-compose.yml}, ubicado en la raíz del repositorio.

El servicio se inicia mediante:

\begin{lstlisting}[language=bash]
docker compose up -d
\end{lstlisting}

El servidor de base de datos queda disponible en el puerto 5432.  
Los scripts de inicialización de la base de datos se encuentran en el directorio \texttt{db/init}, permitiendo la creación automática de tablas y datos de prueba durante la primera ejecución del contenedor.

Para detener el servicio se utiliza:

\begin{lstlisting}[language=bash]
docker compose down
\end{lstlisting}

En caso de requerir una reinicialización completa de la base de datos (incluyendo la eliminación de volúmenes persistentes), se ejecuta:

\begin{lstlisting}[language=bash]
docker compose down -v
\end{lstlisting}

\subsubsection*{Conexión entre Next.js y PostgreSQL}

La comunicación entre la aplicación y la base de datos se realiza mediante el paquete oficial \texttt{pg}, instalado con:

\begin{lstlisting}[language=bash]
npm install pg
\end{lstlisting}

La conexión se gestiona mediante un Pool reutilizable configurado a través de variables de entorno definidas en el archivo \texttt{.env.local}.  

Las consultas a la base de datos se implementan dentro de las API Routes de Next.js, integrando así la capa de acceso a datos dentro del propio proyecto full-stack.

\subsubsection*{Control de versiones}

El control de versiones se configuró mediante Git, inicializando el repositorio local con:

\begin{lstlisting}[language=bash]
git init
git add .
git commit -m "Initial commit"
\end{lstlisting}

Posteriormente, el repositorio fue vinculado a un repositorio remoto alojado en GitHub.

\subsubsection*{Entorno de documentación en \LaTeX}

La memoria se desarrolló en \LaTeX utilizando la distribución MiKTeX junto con Strawberry Perl, necesario para la ejecución de \texttt{latexmk}. La edición y compilación se realizaron en Visual Studio Code mediante la extensión \textit{LaTeX Workshop}.

La compilación del documento se realiza mediante:

\begin{lstlisting}[language=bash]
latexmk -pdf -shell-escape main.tex
\end{lstlisting}

También es posible compilarlo mediante la función integrada en el editor.

\paragraph{Configuración del PATH}

Durante la instalación fue necesario añadir manualmente al \texttt{PATH} del sistema las rutas correspondientes a MiKTeX y Strawberry Perl, ya que inicialmente el comando \texttt{latexmk} no era reconocido por la terminal. Tras actualizar la variable de entorno y reiniciar la sesión, el entorno quedó correctamente configurado.